\scanpage{49}
\begin{flushright}
atkin \mspage{\unsure{9|8|6}}
\end{flushright}

However a significant difference arose here as to the proof of these M- and R-results. If we define for any power series
\[
\handwritten{A(x)=\sum_{n\geqslant N}a(n)x^{n}}
\]
the operator $U=U_{q}$ by
\[
\handwritten{U A(x)=\sum_{qn\geqslant N}a(qn)x^{n},}
\]
then we have
\[
\handwritten{\sum_{n=1}^{\infty}c(q^{\alpha}n)x^{n}=U^{\alpha}j(x).}
\]

As a specific example we now take the simplest case \uline{$q$}$=19$. We can show by modular theory* that there exist functions $g_{n}(x)$ for $n\handwritten{\geqslant}1$
\par\noindent\rule{\linewidth}{0.3pt}\par
\noindent *Of course the same theory could have been applied to the case $q=13$.The main difficulty is that this comparatively deep theory gives for our present purposes results which are \uline{weaker} than those obtained from an elementary use of the modular equation for $q=13$.
\par\noindent\rule{\linewidth}{0.3pt}\par
\noindent with integral coefficients\inserted{ in their power series expansions} and leading terms $x^{n}$,such that
\[
\resizebox{.86\linewidth}{!}{$\displaystyle \handwritten{\begin{aligned}[b]
 U_{19}\,j(x)&\equiv a_{1}\,g_{1}(x)+19\,a_{2}\,g_{2}(x)+19\,a_{3}\,g_{3}(x)\quad(\mathrm{mod}\;19^{2}),\\
 U_{19}\,g_{1}(x)&\equiv c_{11}\,g_{1}(x)+19\,c_{12}\,g_{2}(x)+19\,c_{13}\,g_{3}(x)\quad(\mathrm{mod}\;19^{2}),\\
 U_{19}\,g_{2}(x)&\equiv c_{21}\,g_{1}(x)+c_{22}\,g_{2}(x)+c_{23}\,g_{3}(x)+\\
 &\hspace{13em}+19\,c_{24}\,g_{4}(x)\quad(\mathrm{mod}\;19^{2}),\\
 \text{\deleted{\unsure{etc.,|etc,}}}\quad\vdots&
\end{aligned}}$}
\tag*{\handwritten{(4.1)}}
\]
\reversemarginpar\margin{\handwritten{\tikz[x=1em,y=1em,baseline=-.15em]{\draw (.57,.8)--(.07,-.7); \draw (.44,.51) .. controls (.07,.69) and (.04,.31) .. (.24,.24) .. controls (.02,.18) and (.03,.01) .. (.33,.03);}}}
where the \inserted{\handwritten{$a_{i}$ and}} $\handwritten{c_{ij}}$ are integral constants. Now it turns out \uline{by} \uline{actual} \uline{computation} that some of these constants contain powers of $19$ which were not deducible \uline{a} \uline{prio\mbox{\makebox[0pt][l]{\handwritten{\tikz[x=1em,y=1ex,baseline=0pt]{\draw (.02,-.28)--(.54,-.28)--(.54,1.54)--(1.01,1.6);}}}ir}}; moreover these full powers of $19$ which occur in actual fact are needed to prove the M- and R\normalmarginpar\margin{\handwritten{\tikz[x=1em,y=1em,baseline=-.1em]{\draw (.5,.72)--(.08,-.35) (.27,-.10)--(.38,-.35);} $-$\\(hyphen)}}\deleted{j}\inserted{\handwritten{-}}results. Thus in these cases(\uline{$q$}$=17,19,\text{and }23$)we have at present the situation that the M- and R-results can only be established for a limited number of
% archive sheet number: 48
